\section{引言}
\label{sec:intro}

告警触发后，分析师要判断它的起因：正在进行的攻击、配置错误、经授权的活动，还是误报。分析师通过查询认证日志、请求模式、配置、变更工单和威胁情报来缩小候选解释的范围；当某条观测支持一种解释、其余解释都被排除时，就停止调查。结论决定该事件是否升级处理，而每一次查询都要花时间。

手工做这件事无法扩展。告警数量超出分析师的调查能力，而且大多数告警最终是良性的~\citep{nodoze}。能自己选择查询、阅读结果的 LLM 智能体天然契合这个循环。然而，许多组织不能或不愿把原始遥测数据发送给托管模型。对它们来说，现实的选择是一个小到能在单块本地 GPU 上运行的开源权重模型；而近期关于 LLM 驱动的攻击调查的工作已经报告，这类模型在复杂案例上落后于闭源模型~\citep{clouseau}。

自动化分诊和调查的工作大体沿两条路线展开。第一条构建主机活动的溯源图，并用固定规则或学习模型对其排序或摘要~\citep{nodoze,kairos}。这类系统减少了分析师需要阅读的内容，但不决定下一步运行哪个查询。第二条使用 LLM 智能体来选择查询并解释结果~\citep{excytin,clouseau}，采用 ReAct 模式，即每一步都由模型决定~\citep{react}。这些智能体把调查流程交给模型，而对小型本地模型来说，失败的恰恰是流程。在我们的预注册研究中，自己选择探针的 Qwen2.5-7B 只在 8\% 到 13\% 的分诊案例中以已验证的结论结束：它一直探测直到预算耗尽。Llama-3.1-8B 从不下结论。在 3{,}389 次决策中，它每次都要求再多一个探针，即使证据已压倒性充分、也已没有合法探针可用。更大的模型会下结论，但在没有引导的情况下，其引用中有 22\% 到 31\% 是无效证据，并把 25\% 到 37\% 的需处置事件判为良性。

本文提出 \hesp{}，一个把调查流程放在本地 LLM 智能体之外的控制器。一次调查由两个决策组成：\emph{下一步探什么}，以及\emph{证据何时已足够}；控制器可以接管其中之一或两者都接管。\hesp{} 维护一个记录相互竞争的解释的账本，用表格按贝叶斯规则更新；这些表格说明每个探针在每种解释下会返回什么。它按每单位成本的期望信息增益为只读探针排序，只接受有当前支持性观测的结论，并能自行结束调查。模型仍然是规划器：它读取账本，可以提议探针、给出结论或放弃。每条预测都在其观测之前写入日志，每个回合事后都会被审计。

我们用来自两个家族的五个开源权重模型（7B 到 72B）做了四项预注册研究，在受控的分诊环境中共 7{,}272 个经过审计的回合。从无 LLM 的开发回合中计数得到的表，把 Qwen2.5-7B 的已验证完成率从 0.125 提升到 1.000，与从环境生成器直接得到的表持平。在规划器所获信息保持不变的情况下，信息增益排序在每个会下结论的模型上带来 $+0.26$ 到 $+0.35$ 的提升。控制器侧停止把 Llama-3.1-8B 从 0 提升到 0.917；在此之后，控制器选探针对它不再带来可检测的完成率提升，而 Qwen2.5-7B 两部分都需要。完全不用 LLM 的控制器本身完成 0.917 的案例；最好的 LLM 配置达到 1.000，因为强模型会等到确认性证据出现才下结论。因此在我们的环境中，收益来自把流程从模型手中拿走，而不是模型调查得更好：每个起因都会产生一个独特的离散观测，环境自己把响应翻译成这些观测，似然表和验证器也都来自同一个生成器。我们据此解读结果：它说明的是“谁应该掌控流程”，至于模型在观测模糊时能否有所贡献，本文没有回答。

概括而言，本文的贡献如下：
\begin{itemize}
  \item 我们表明，小型本地 LLM 在告警分诊上的失败发生在流程层面，并且“选择探什么”和“决定何时停”是两种独立的失败，在规模相近的模型之间也各不相同（\S\ref{sec:motivation}，\S\ref{sec:rq-stop}）。
  \item 我们设计、实现并开源了 \hesp{}\footnote{代码、冻结的协议、预测表、结果文件和回合日志见 \url{https://github.com/lzwhehe/HESP}。}，一个为本地 LLM 规划器接管假设账本、探针选择、结论接受和停止的控制器，带有预写日志和逐回合审计（\S\ref{sec:approach}）。
  \item 我们用五个开源权重模型做了四项预注册研究，采用冻结的协议和源码哈希、在不同选择器之间保持规划器提示不变的盲对照，以及把探针选择与停止分开的析因设计（\S\ref{sec:setup}，\S\ref{sec:results}）。
  \item 我们表明，在这个环境中，完全不含 LLM 的控制器完成 0.917 的案例，与放入控制器的五个模型中的四个持平，因此可测的收益来自流程本身；剩余的差距来自强模型等待确认证据（\S\ref{sec:rq-stop}）。从同一生成器的少量无 LLM 运行中计数得到的似然表，与生成器自身的表一致，而由模型给出的表则不行（\S\ref{sec:rq-help}）；由于计数和评估共用生成器，这只说明似然表不必手写，不说明它能迁移到真实数据。
  \item 我们测量了三个公开安全数据集，说明为什么它们都无法检验调查策略，并给出一个基于真实攻击遥测记录的案例研究，展示这些数据集所缺少的那类案例（\S\ref{sec:realdata}）。
\end{itemize}
